\subsection{The exact theta-orbit criterion and its remaining arithmetic condition}

\begin{claim}{A.C2.5}{Proved}
Let $C$ be a smooth proper genus-two curve over an algebraically closed
field $\Omega$ of characteristic two. Let $\mathcal B$ be its line bundle
of locally exact differentials, with the absolute-Frobenius module
structure, and put
\[
 \Theta=\{L\in\operatorname{Pic}^0(C):
                   H^0(C,\mathcal B\otimes L)\ne0\}.
\]
For an odd integer $n>1$, let $L$ have exact order $n$, and let
$D\to C$ be the associated connected cyclic \'etale cover. Write
$\mathcal O_n$ for the orbits of multiplication by two on
$(\mathbf Z/n)\setminus\{0\}$. Then
\[
 f(D)=f(C)+
 \sum_{\substack{O\in\mathcal O_n\\
                  L^i\notin\Theta\ \text{for all }i\in O}}|O|.
                                                               \tag{C2.16}
\]
The set $\Theta$ is symmetric under $L\mapsto L^{-1}$.
If every orbit meets $\{i:L^i\in\Theta\}$, a lift of the
hyperelliptic involution has a quotient $B$ of genus $(n-1)/2$
and $2$-rank zero. Two deck-conjugate quotient maps $D\to B$
jointly separate geometric points. In particular, for an odd prime
$\ell$, the sufficient condition is satisfied whenever
\[
 L\in\Theta[\ell]\setminus\{0\},\qquad
               \langle2,-1\rangle=(\mathbf Z/\ell)^\times.       \tag{C2.17}
\]
Over a field finitely generated over $\F_2$, a curve satisfying
this geometric condition has an injective geometrically
prime-to-two point-section map, as do its finite-extension forms
and hyperbolic open subcurves.
\end{claim}
\begin{proof}
The exact sequence (C2.15) is valid for every smooth curve in
characteristic two. Its local description also proves
$\mathcal B^{\otimes2}\simeq\Omega_C^1$: if $t$ is a local
parameter, a locally exact differential is $h(t)^2dt$, so the
transition functions for $\mathcal B$ square to those for
$\Omega_C^1$. Thus $\deg\mathcal B=1$. Riemann--Roch gives
\[
 h^0(\mathcal B\otimes M)
   =h^0(\Omega_C^1\otimes\mathcal B^{-1}\otimes M^{-1})
   =h^0(\mathcal B\otimes M^{-1})
\]
for every degree-zero $M$, proving the symmetry.

For $0\ne i\in\mathbf Z/n$, both $L^i$ and $L^{2i}$ are
nontrivial of degree zero. They have no global sections, and
their first cohomology spaces have dimension one. Tensoring
(C2.15) by $L^i$ therefore identifies the kernel of
\[
 F_i:H^1(C,L^i)\longrightarrow H^1(C,L^{2i})
\]
with $H^0(C,\mathcal B\otimes L^i)$. Hence $F_i$ is zero
precisely when $L^i\in\Theta$, and otherwise it is a bijective
semilinear map between one-dimensional spaces. The cyclic-cover
decomposition
\[
 H^1(D,\mathcal O_D)=
             \bigoplus_{i\in\mathbf Z/n}H^1(C,L^i)
\]
is respected by Frobenius, with indices multiplied by two.
On an orbit meeting $\Theta$, the $|O|$th Frobenius power is
zero. On an orbit disjoint from $\Theta$, Frobenius is bijective
and its stable rank is $|O|$. The index-zero summand has stable
rank $f(C)$. The Artin--Schreier description of $2$-rank now
proves (C2.16). This proof treats all nonzero characters,
including characters of proper conductor when $n$ is composite.

The hyperelliptic involution $\iota$ acts by $-1$ on
$\operatorname{Pic}^0(C)$, so it lifts to an automorphism
$\tau$ of the cyclic cover which conjugates a deck generator
$\delta$ to $\delta^{-1}$. Such a lift is an involution:
$\tau^2\in\langle\delta\rangle$ commutes with $\tau$, and
inversion fixes no nonidentity element of an odd cyclic group.
At a branch point of $C\to\mathbf P^1$, the action of $\tau$
on the $n$ points above it is an affine inversion of the cyclic
torsor, so exactly one of those points is fixed. The completed
local map $D\to C$ at that point is an isomorphism, and the
different exponent for $D\to D/\langle\tau\rangle$ equals
the exponent for $C\to\mathbf P^1$. The latter different has
total degree six, since $g(C)=2$. Since $g(D)=n+1$,
Riemann--Hurwitz gives
\[
       2n=2(2g(B)-2)+6,\qquad g(B)=(n-1)/2.
\]

The dihedral argument of A.C2.2 applies to every odd $n$:
\[
                         J_D\sim J_C\times J_B^2.
                                                               \tag{C2.18}
\]
Here is the precise linear-algebra verification. In an auxiliary
rational Tate module at a prime different from two and from the
divisors of $n$, let $W$ be the complement of the cyclic
invariants. Its nontrivial cyclic characters occur in pairs
$\lambda,\lambda^{-1}$, and $\tau$ exchanges the members of
each pair. Consequently $W^\tau$ has half the dimension of $W$.
The subspaces $W^\tau$ and $\delta W^\tau$ intersect trivially:
an element in the intersection is fixed by both corresponding
reflections, hence by $\delta^2$, hence by $\delta$, whereas
$W$ has no cyclic invariants. Both subspaces are pulled-back
Tate modules of the quotient Jacobian. They therefore span $W$
and give (C2.18). The cyclic-invariant part is $VJ_C$, on which
$\tau$ acts by $-1$.

If every orbit meets $\Theta$, formula (C2.16) and additivity
of $2$-rank under isogeny give $f(B)=0$. Two quotient maps
associated with the reflections $\tau$ and
$\delta^{-1}\tau\delta$ jointly separate points: a nontrivial
collision would be a fixed point of $\delta^2$, impossible
because the odd cyclic deck action is free. Composing with
Abel embeddings into $J_B$ and applying A.10 proves the
point-section assertion. The covers, quotient maps and chosen
base points descend after a finite extension, as required there.

Finally suppose $n=\ell$ is prime and (C2.17) holds. If
$L\in\Theta$, symmetry also gives $L^{-1}\in\Theta$.
Every nonzero residue modulo $\ell$ then belongs to the
doubling orbit of $1$ or of $-1$, proving the required
orbit condition. In particular, the condition is weaker than
requiring two itself to be a primitive root: for $\ell=7$
there are two doubling orbits, exchanged by sign, and both
are met simultaneously.
\end{proof}

\begin{claim}{A.C2.6}{Unresolved}
For an arbitrary genus-two curve over $\overline{\F}_2$, the
existence of an odd-order torsion point satisfying the
all-orbits condition of A.C2.5 is not established here.
Prime-to-two torsion density on $\Theta$ does not supply the
missing order and orbit restrictions.
\end{claim}
\begin{proof}[Exact remaining obligation]
The theta curve is a translate of the Abel image of $C$ and
generates its Jacobian. Yang's Proposition 2.3 therefore gives
density of odd-order torsion on it over $\overline{\F}_2$
\citep{A_Yang2016}. That statement does not prescribe a prime
order $\ell$, the subgroup $\langle2,-1\rangle$ modulo $\ell$,
or which other multiples of a chosen point remain on $\Theta$.
Formula (C2.16) identifies the exact missing condition: every
nonzero doubling orbit must contain a theta point. For a
composite order, the orbits of proper-order characters must
also be treated. A single point of exact order $n$ only
guarantees a zero map on its own orbit and, by symmetry,
on the opposite orbit.

This distinction is substantive. In the $2$-rank-one genus-two
family of A.C2.2 every cyclic cubic \'etale cover has an ordinary
elliptic involution quotient. Formula (C2.16) therefore implies
$\Theta\cap J_C[3]=\{0\}$ in this family. Thus even the
smallest prescribed odd prime need not occur on the nonzero
theta locus. Any cyclic cover whose degree is divisible by
three retains an ordinary contribution from these order-three
characters, regardless of theta intersections among its
primitive characters. The ordinary family off the hypersurface
of A.C2.1 has the same nonzero-three-torsion obstruction.

Nor does a single theta intersection give a quantitative
small-$2$-rank tower. Formula (C2.16) gives
\[
 \frac{f(D)}n=\frac{f(C)}n+
   \frac1n\sum_{O\text{ disjoint from }\Theta}|O|.
\]
To make this ratio small, theta intersections must eliminate
almost all characters. Density alone supplies no such
proportion. No assertion that arbitrary curves admit these
covers, or that normalized $2$-rank tends to zero in suitable
prime-to-two towers, is used in the positive conclusions.
\end{proof}
